% =====================================================================
%  2 -- SILICON BURNING.
%
%  Source: tier3.md 0.3 intro (871-876), 0.3.1 (878-943),
%          0.3.2 (945-1024), 0.3.3 (1026-1064), 0.3.4 (1066-1121),
%          0.3.5 (1123-1170), 0.3.6 (1172-1259);
%          plus self-check Q13 from 0.10 (2746-2749) -> sec:selfcheck-siburn.
%  Headings covered: "0.3 Silicon burning proper" (-> \section),
%          "0.3.1 The mechanism, mechanically",
%          "0.3.2 The two clusters, and what the group boundary means
%                 physically",
%          "0.3.3 The endpoint, and why Ye chooses it",
%          "0.3.4 Hydrostatic versus explosive silicon burning",
%          "0.3.5 alpha-rich freeze-out -- the third regime",
%          "0.3.6 The NSE transition -- where a network is replaced by a
%                 table".
%  Re-homings / rewordings: none of substance.  Internal cross-refs
%          (0.3.3, 0.2.2, IV.3, 0.6, 0.6.4, 0.7.3, 0.2.3, 0.3.5, III,
%          0.8) are live \ref's per the registry; 0.7.3 and III/IV.3 are
%          forward references in this edition, which the source's own
%          wording already tolerates ("SSIV.3 is about that gap").
%          The 0.3.4 [assumed band -- ...] tag is rendered as a one-off
%          \provtag (the \assumednote macro would print an extra dash).
% =====================================================================

\section{Silicon burning}
\label{part:siburn}

Tier~2 \S 0.2 gives the mechanism in outline. This section is the fuller
physical picture, including the part Tier~2 explicitly deferred: silicon
burning happens in \textbf{two different astrophysical settings}, with different
consequences, and the project's data is a model of only one of them.

% =====================================================================
\subsection{The mechanism, mechanically}
\label{sec:mechanism}

Start from a zone of \nuc{Si}{28} and \nuc{S}{32} (the ash of oxygen burning;
\citealt[\S 1]{hix1996}: ``silicon isotopes and their $\alpha$-nuclei
neighbors, S, Ar, and Ca'') at $\Tn \approx 3.5$, $\rho \approx 10^{8}$.

\textbf{Step 1 --- the photon bath finds the weakest link.} The Planck
distribution at $kT = 0.30$~MeV has an exponentially small but nonzero
population above 10~MeV (\citealt{hix1996}: a photodisintegration channel
becomes important once its Q-value falls below $\approx 30\,\kB T$). Nuclei
with the \emph{lowest} particle-separation energies are attacked first: the
($\gamma$,$\alpha$), ($\gamma$,p) and ($\gamma$,n) channels on whichever
species in the mixture happens to be least bound against that channel. For a
silicon-group mixture, that is not \nuc{Si}{28} --- \nuc{Si}{28} is
comparatively tightly bound ($S_\alpha = 9.98$~MeV, $S_p = 11.6$~MeV, the
highest along the $\alpha$-chain; AME2020, \citealt{wang2021}) --- but the
odd-$A$ and less-bound neighbours.

\textbf{Step 2 --- a reservoir of light particles appears.} Every
photodisintegration deposits a free $\alpha$, p or n into the plasma. These
particles have no Coulomb barrier problem worth speaking of (n) or a small one
(p, $\alpha$ on light targets), so their capture rates are enormous. A free
neutron at these conditions is captured in a time of order $10^{-10}$~s ---
and that figure is deliberately conservative:
$n_{\mathrm{target}}\langle\sigma v\rangle$ with
$n_{\mathrm{target}} \sim 10^{30}$~cm$^{-3}$ and
$\langle\sigma v\rangle \sim 10^{-17}$~cm$^{3}$\,s$^{-1}$ gives closer to
$10^{-13}$~s, so the true separation from the bridge timescales of step 4 is
wider than the number quoted here. \derivedhere[order of magnitude]

\textbf{Step 3 --- the reservoir is spent on whatever is most bound.} The
captured light particles are distributed over the whole silicon group by the
fast capture/photodisintegration pairs, and net flow proceeds toward higher
binding energy. The material walks up in $A$, one light particle at a time.

\textbf{Step 4 --- the flow encounters bottlenecks.} Not every step is fast. A
few reactions have small rates relative to the rest --- because of an
unfavourable Q-value, a Coulomb barrier on a heavier target, or a structural
feature of the nucleus. Those are the \textbf{bridges}: the reactions the net
flow has to pass through, connecting the silicon group to the iron group.
Everything else is running back and forth in near-balance.

\textbf{Step 5 --- the endpoint is set by \Ye.} Material accumulates in the
iron peak, and \emph{which} iron-peak nuclei depends on the neutron-to-proton
ratio available, i.e. on \Ye. \S\ref{sec:endpoint}.

The end-to-end result: $2\,\nuc{Si}{28} \to \nuc{Ni}{56}$, releasing
0.195~MeV/nucleon (\S\ref{sec:energypergram}; $Q = 10.922$~MeV from the AME2020
mass excesses), achieved by disassembling some fraction of the silicon
completely and reassembling it (\citealt{hix1996}: ``a complex series of
photodissociation and capture reactions''). No \nuc{Si}{28} nucleus ever
touches another \nuc{Si}{28} nucleus.

\textbf{The signature of this mechanism in the equations} --- the reason it
dominates everything in Tiers~1--2 --- is that steps 1--3 are \emph{fast and
nearly balanced}, while step 4 is \emph{slow and one-directional}. Fast
near-balanced pairs are exactly:

\begin{itemize}
\item \textbf{near-zero \kap} (Tier~2 \S II.4): $|\netflux| \ll \fp + \fm$;
\item \textbf{large Jacobian eigenvalues with tiny net effect}: stiffness
  (Tier~2 \S IV.1);
\item \textbf{quasi-statistical equilibrium}: the balanced subnetwork's
  composition is a function of a few parameters (Tier~2 \S 0.1);
\item \textbf{catastrophic cancellation in floating point} (Tier~0 \S V.1).
\end{itemize}

Those are four names for step 1--3. And the \emph{interesting} physics --- the
net flow, the energy release, the \Ye{} evolution --- lives entirely in step 4
and in the weak sector, which is a small fraction of the columns
(\citealt[\S 6]{hix1996}: one can ``ignore the reactions within the
quasi-equilibrium groups \ldots\ [and] concentrate on those reactions which
enter or leave the QSE groups''). That asymmetry is the project's central bet.

% =====================================================================
\subsection{The two clusters, and what the group boundary means physically}
\label{sec:twoclusters}

The balanced subnetwork does not connect everything to everything. It splits:
a \textbf{silicon group} (roughly $24 \le A \lesssim 45$) and an \textbf{iron
group} ($A \gtrsim 45$), each internally equilibrated on a timescale short
compared with the burning, but connected to each other only through the slow
bridges (\citealt[\S\S 4, 6.1]{hix1996}: ``two groups, roughly divided by
$A \simeq 45$'', the bottom edge of the silicon group being Mg).
Quasi-equilibrium itself is \citeauthor{bodansky1968}'s
(\citeyear{bodansky1968}) result --- a \emph{single} group from \nuc{Si}{28} to
the iron peak; the split into two groups joined by bridges is
\citeauthor{woosley1973}'s (\citeyear{woosley1973}), confirmed and made
\Ye-dependent by \citet{hix1996} (an earlier version of this sentence credited
the two-group split to BCF68; corrected in the 2026-08-18 audit). It is the
reason the composition during silicon burning is neither a free ODE trajectory
(too many degrees of freedom) nor a state function (NSE --- too few).

\textbf{Why the split exists at all} is worth understanding physically rather
than taking as an empirical fact. The equilibration of a group requires that
every member be connected to every other by reactions fast compared with the
burning timescale. In the $A \approx 45$ region there is a gap in the reaction
network's effective conductivity, and it has two causes pulling the same way.

\textbf{(a) The boundary nuclei are \emph{loosely} bound against $\alpha$
emission, not tightly.} This is the one that is easy to get backwards. The
$\alpha$-separation energies say:
%
\begin{equation}
S_\alpha(\nuc{Ca}{40}) = 7.04\ \mathrm{MeV},
\qquad
S_\alpha(\nuc{Ti}{44}) = 5.13\ \mathrm{MeV},
\label{eq:salpha}
\end{equation}
%
\sourcednote{AME2020 via the IAEA Live Chart --- $Q_\alpha(\nuc{Ca}{40}) =
-7039.78$~keV, $Q_\alpha(\nuc{Ti}{44}) = -5127.10$~keV} --- so \nuc{Ti}{44}
photodisintegrates \emph{more} readily than the \nuc{Ca}{40} below it, and the
($\gamma$,$\alpha$) reverse on the \nuc{Ca}{40}($\alpha$,$\gamma$)\nuc{Ti}{44}
step runs hard. The bridge is slow in \textbf{net} terms precisely because its
reverse is fast: material pushed across is pushed straight back. A local dip in
the binding-energy-per-nucleon curve (\nuc{Ti}{44} at 8.534~MeV/nucleon,
\emph{below} \nuc{Ca}{40}'s 8.551) is what a bottleneck looks like from the
thermodynamic side. (This is the $\alpha$-chain step
\nuc{Ca}{40}($\alpha$,$\gamma$)\nuc{Ti}{44}; the (p,$\gamma$) bridges into
\nuc{Ti}{46} feed a tightly bound product, $S_p(\nuc{Ti}{46}) = 10.3$~MeV, and
are limited instead by the small feeder abundances ---
\citealt[\S 6.2]{hix1996}.)

\textbf{(b) The capture reactions that would carry flow across face Coulomb
barriers on targets with $Z \approx 20$--21}, which suppresses the forward
direction independently of (a) (\citealt[\S 1]{hix1996}, after
\citealt{thielemann1985}: the bottleneck ``coinciding roughly with $Z = 21$'' is
bridged on the proton-rich side).

Together these make the crossing slow relative to the silicon-group internal
reactions but not slow enough to stop the flow entirely --- which is exactly the
condition for two separately-equilibrated groups joined by identifiable bridges.

\textbf{And therefore the boundary is not sharp} (\citealt[\S 4]{hix1996}:
``some ambiguity in defining the boundary between the QSE groups''). Whether
\nuc{Sc}{45} counts as silicon-group or iron-group is a \emph{modelling choice
about a continuum}. The repo handles this correctly by making it a config
variant (\code{configs/qse\_groups.yaml}: \code{default} $24 \le A < 45$ vs
\code{a24\_46} with the boundary at $A = 46$, putting \nuc{Sc}{45} inside the
group) and requiring every group/bridge analysis to report both
(\code{CLAUDE.md}). The measured consequence --- the verdict is insensitive to
the variant at the $\le 0.02$ level in top-$k$ shares,
\resultsref{2026-07-12} --- is exactly the kind of result the discipline exists
to produce: the choice was declared a choice \emph{before} it turned out not to
matter.

\textbf{The bridge sets, measured.} For \msn{80} on the relaxed manifold:
Ne22($\alpha$,n)Mg25 (share 0.20), Al27(p,$\alpha$)Mg24 (0.13),
P31(p,$\alpha$)Si28 (0.11), Na23($\alpha$,p)Mg26, Mg24(n,$\gamma$)Mg25; top-20
columns carry 0.88 of the inter-group flux. For \msn{151}: Mg24(p,$\alpha$)Na21
(0.14), Mg24(p,$\gamma$)Al25, Cl35(p,$\alpha$)S32, P31(p,$\alpha$)Si28,
Ca44(p,$\gamma$)Sc45 --- and the literature bottleneck
\textbf{\nuc{Sc}{45}(p,$\gamma$)\nuc{Ti}{46} at rank 2 of 251} under the
\code{a24\_46} boundary. \resultsref{2026-07-12}

That last one deserves a moment. The reaction was named in the literature as
the high-\Ye{} bottleneck of silicon burning (\citealt[\S VIb]{woosley1973};
\citealt[\S 6.2]{hix1996}: \nuc{Sc}{45}(p,$\gamma$)\nuc{Ti}{46} ``accounts for
approximately 70\% of the entire flux into the iron peak group'' at
$\Ye = 0.498$ --- and \emph{not} the link at $\Ye = 0.46$; an earlier version of
this sentence cited \citet{grichener2025} for the claim, whose paper does not
mention the reaction; corrected in the 2026-08-18 audit), and this project
found it independently at rank 2 of 251 inter-group carriers on relaxed
high-\Ye{} trajectory rows in the 3.3--5~GK window. That is a genuine external
validation of the whole flux machinery: a published physical claim, reproduced
by an independent pipeline (one that shares the REACLIB/pynucastro rate-library
family with the literature, so ``independent'' is about the code, not the
inputs). It is also a warning about distribution --- on the \emph{training}
distribution (random Sobol compositions) the same reaction ranks only 31/251
\resultsref{2026-07-10}. The bottleneck is a property of relaxed flow, and it
is invisible where the model will be trained. \S\ref{sec:twodistributions} is
about that gap.

% =====================================================================
\subsection{The endpoint, and why \texorpdfstring{\Ye}{Ye} chooses it}
\label{sec:endpoint}

Silicon burning ends when the material reaches the iron peak and there is
nothing more exothermic to do. \emph{Which} iron-peak nucleus is the endpoint
is a selection problem, and the selection rule is \Ye.

Consider NSE at fixed $(T, \rho, \Ye)$. The composition minimises free energy
subject to $\sum X_i = 1$ and $\sum (Z_i/A_i)X_i = \Ye$ (the chemical-potential
formulation of NSE, \citealt[\S 5]{hix1999b}, after \citealt{clifford1965}). At
low enough temperature that the binding-energy term dominates the entropy term,
the winner is the nucleus that maximises binding energy per nucleon
\textbf{subject to having the right $Z/A$} (\citealt{hix1999b}: ``the most
abundant nuclei at low temperatures are the most bound nuclei for which
$Z/A \sim \Ye$''; \citealt[Fig.~2]{hartmann1985}). Since $\Ye = \langle Z/A
\rangle$ for the mixture, and a single dominant species must therefore have
$Z/A \approx \Ye$:

\begin{center}
\begingroup\small
\begin{tabularx}{\textwidth}{@{}L{1.4cm} L{3.4cm} Y@{}}
\toprule
\textbf{\Ye} & \textbf{$Z/A$ match} & \textbf{dominant NSE species} \\
\midrule
0.500 & $28/56 = 0.500$ & \textbf{\nuc{Ni}{56}} (doubly magic, $B/A = 8.643$) \\
0.482 & $26/54 = 0.4815$ & \textbf{\nuc{Fe}{54}} (and \nuc{Ni}{58},
  $28/58 = 0.483$, $B/A$ 8.732 vs 8.736) \\
0.464 & $26/56 = 0.4643$ & \textbf{\nuc{Fe}{56}} ($B/A = 8.790$) \\
0.452 & $28/62 = 0.4516$ & \nuc{Ni}{62} region (\nuc{Fe}{58}, at
  $26/58 = 0.448$, sits just below) \provtag{[dominance assumed here; no
  source opened]} \\
\bottomrule
\end{tabularx}
\endgroup
\end{center}

\derivedhere[{the $Z/A$ arithmetic recomputed, $B/A$ values AME2020; the NSE
dominance rule is \citealt[\S 5]{hix1999b} and \citealt[Fig.~2]{hartmann1985},
and the repo's own Saha solver reproduces it: \resultsref{2026-07-10} NSE fe56
$= 0.69$ at $\Ye = 0.467$ and 2026-07-11 fe56 $= 0.72$ at $\Ye = 0.472$.}]

So the composition arriving at the iron peak is a \textbf{direct readout of
\Ye}, and the \nuc{Ni}{56} mass --- the thing that powers the supernova light
curve \citep{boccioli2024, arnett1982} --- is set by how much of the burnt
material stayed near $\Ye = 0.5$.

\textbf{This is the reason the entire project is organised around \Ye{} and not
around \nuc{Ni}{56}.} \Ye{} is upstream; \nuc{Ni}{56} is downstream. An
emulator that gets \Ye{} right gets the endpoint right for the right reason.
One tuned to reproduce \nuc{Ni}{56} can be right about \nuc{Ni}{56} and wrong
about the physics --- and will then be wrong about \nuc{Ni}{57}, \nuc{Ni}{58},
\nuc{Ti}{44} and everything else at once, because those \emph{also} depend on
\Ye{} and the tuning did not.

% =====================================================================
\subsection{Hydrostatic versus explosive silicon burning}
\label{sec:hydroexplosive}

Tier~2 deferred this and it matters, because the two settings are physically
different and the project's data models one of them.

\textbf{Hydrostatic (pre-supernova) silicon burning} is what
\S\ref{sec:neutrinoclock} described: the core burns for $\sim$a day (0.01~yr
for a 20~\Msun{} star in \citealt{hix1999b}'s Table~1) at
$\Tn \approx 3$--4, $\rho \approx 10^{7}$--$10^{9}$ (\citealt{hix1996} study
3.5--5~GK and $10^{7}$--$10^{10}$~\gcc{} as the core-Si-burning cube), cooling
by neutrinos, in quasi-equilibrium, with electron captures having plenty of
time to act. \Ye{} falls substantially. The product is the iron core.

\textbf{Explosive silicon burning} happens seconds later, when the bounce shock
passes through the silicon and oxygen shells. Conditions: peak
$\Tn \approx 4$--5+ (\citealt{boccioli2024}: incomplete Si burning for
$4 < \Tn < 5$, full NSE above 5; \citealt{hix1999a}: $T_{9i} < 5$ incomplete),
$\rho \approx 10^{7}$--$10^{8}$, and a duration set by the \emph{hydrodynamic
expansion timescale} --- $\tau_{\mathrm{HD}} = (24\pi G\rho)^{-1/2} =
446\,\rho_6^{-1/2}$~ms (\citealt{fowler1964}, as used by \citealt[eq.~1]{hix1999a}),
i.e. 0.05--0.4~s for $\rho = 10^{6}$--$10^{8}$~\gcc, with the temperature
falling on $3\tau_{\mathrm{HD}}$. The material is heated, burns, and then
expands and cools on that same timescale.

Three differences with real consequences:

\begin{center}
\begingroup\small
\begin{tabularx}{\textwidth}{@{}L{3.0cm} Y Y@{}}
\toprule
 & \textbf{hydrostatic} & \textbf{explosive} \\
\midrule
duration at peak $T$ & $\sim 10^{5}$~s & $\sim 10^{-1}$~s \\
weak reactions & have time; \Ye{} drops & \textbf{frozen}: $\Ye \approx$
  inherited value \\
endpoint & true QSE $\to$ NSE approach & interrupted; freeze-out matters \\
product & iron core (collapses) & \textbf{the ejecta} (observable) \\
\bottomrule
\end{tabularx}
\endgroup
\end{center}

\textbf{The weak-freeze difference is the important one}
(\citealt[\S 2]{hix1996}: ``the timescales for changes in \Ye{} via weak
reactions \ldots\ are much longer than the timescales for the strong and
electromagnetic reactions'', so they treat \Ye{} as a constant parameter).
Electron-capture rates across the box span $10^{-5}$--$10^{2}$~s$^{-1}$ per
nucleus (Tier~1 \S V.1 table; magnitudes from \citealt[Figs.~5--6]{langanke2000})
\provtag{[\textsc{assumed} band --- the upper decades belong to collapse
densities $\rho\Ye \gtrsim 10^{9}$--$10^{10}$; retire by tabulating
$\lambda_{\mathrm{EC}}$ for the \S\ref{sec:yecarriers} top-20 at the box
corners; an earlier version quoted $10^{-2}$--$10^{0}$~s$^{-1}$, inconsistent
with Tier~1's own table]}. Over $10^{5}$~s that changes \Ye{} by a lot; over
0.1~s it changes it by $\sim 10^{-2}$ at most and often much less. So
\textbf{explosive burning largely inherits the \Ye{} that hydrostatic burning
produced} --- the \Ye{} statement itself is the HT96 timescale argument above;
the shock's \emph{thermodynamic} conditions likewise ``are determined by the
pre-explosive structure'' \citep{woosley2007} --- and the ejecta composition
is a fossil of the pre-supernova core structure.

Which closes the loop on why the project's regime box is the right one: the box
is the hydrostatic-burning box, and hydrostatic burning is what \emph{sets} the
\Ye{} that everything downstream inherits. The explosive episode redistributes
it into species but does not create it.

\textbf{A caveat the study notes should carry.} The box (\Tn{} 1.6--7.9, $\rho$
$10^{7}$--$10^{9}$) is wide enough to overlap explosive conditions too, and the
training data is a Sobol grid over that box with random compositions --- which
corresponds to no physical setting in particular. Whether the induced sampling
measure resembles either the hydrostatic or the explosive trajectory density
has never been checked; that is \foreignreg{2}{R-07}, and
\S\ref{sec:twodistributions} shows it is not an academic question.

% =====================================================================
\subsection{\texorpdfstring{$\alpha$}{alpha}-rich freeze-out --- the third regime}
\label{sec:alpharich}

One more setting, because it produces two of the observables in
\S\ref{part:observables} and it is the case where ``NSE'' most badly misleads.

Material shocked to \textbf{high temperature and low density} ($\Tn \gtrsim 5$
with $\rho \lesssim 10^{7}$--$10^{8}$~\gcc{} --- equivalently high radiation
entropy; the boundary between normal and $\alpha$-rich freeze-out is a line of
constant entropy per gram, with $\rho_i \approx 10^{8}$ separating the two in
\citealt[\S 2 and Fig.~1]{hix1999a}; the canonical parameter study is
$\Tn = 5.5$, $\rho = 10^{7}$, $s/k \approx 5$, \citealt{the1998}) reaches (or
nearly reaches) NSE, which at those conditions is dominated by free $\alpha$
particles and nucleons rather than by the iron peak
(\citealt[\S 5]{hix1999b}) --- high temperature favours the high-entropy,
many-particle state. As the material then expands and cools, the composition
tries to follow NSE back down toward the iron peak. Reassembling $\alpha$
particles into heavy nuclei requires the \textbf{triple-$\alpha$ reaction} (and
\nuc{He}{4}($\alpha$n,$\gamma$)\nuc{Be}{9}), which is a three-body process and
therefore has a rate $\propto \rho^{2}$ (\citealt{hix1999a}: ``the quadratic
density dependence of the rates for $\alpha+\alpha+\alpha \to \nuc{C}{12}$ and
$\alpha+\alpha+\mathrm{n} \to \nuc{Be}{9}$''; \citealt{the1998}). At low density
that rate cannot keep up with the expansion.

The result is a \textbf{freeze-out with leftover $\alpha$ particles} --- 1--10\%
by mass (the \nuc{He}{4} contours of \citealt[Fig.~1]{hix1999a}; 8\% in their
$\Tn = 6$, $\rho = 10^{7}$ example) --- which then capture onto the iron-peak
nuclei that did form. Consequences:

\begin{itemize}
\item extra \nuc{Ti}{44} (an $\alpha$-capture product; observable,
  \S\ref{sec:remnants}; \citealt{the1998}; \citealt{jordan2003});
\item shifted iron-peak isotopic ratios relative to a normal freeze-out
  (\citealt[\S 5]{hix1999a});
\item a composition that is \emph{not} NSE and, in the canonical $\alpha$-rich
  freeze-out example (pure \nuc{Si}{28} at $\Tn = 5.5$, $\rho = 10^{7}$;
  \citealt{the1998}, after \citealt{meyer1998}: a QSE with a heavy-nucleus
  deficit), never fully was, despite having passed through NSE conditions.
\end{itemize}

\textbf{Why this belongs in a study document about a network emulator.} It is
the cleanest example of the general rule that \textbf{passing through
equilibrium conditions does not make the composition an equilibrium function}
--- the history matters, because the \emph{rate} of departure from equilibrium
competes with the rate of change of the conditions. That is the same statement
as ``QSE is a reduction, not a solution'' (Tier~2 \S 0.1.1), and the same
statement as ``NSE-as-a-predictor is a baseline that has never been run''
(\foreignreg{2}{R-16}). It is also, incidentally, the physical reason that the
low-density corner of this project's regime box is the awkward one --- Tier~1
\S VIII.C.6's $e^{+}e^{-}$ pair problem lives at exactly ($\Tn = 7.9$,
$\rho = 10^{7}$).

% =====================================================================
\subsection{The NSE transition --- where a network is replaced by a table}
\label{sec:nsetransition}

There is a boundary inside the regime box that no tier has treated, and it is
the boundary the emulator's \emph{domain} should probably be defined against.

\textbf{Some stellar-evolution and supernova hydro codes stop integrating the
network at high temperature --- MESA does not.} Above a code-parameter
threshold \typical[{$\sim$5--7~GK; \citealt[\S 5.2]{paxton2015} quote the
generic switch ranges as $> 3$~GK (QSE) and $> 5$~GK (NSE); KEPLER's actual
parameter not read here}] the strong reactions are so fast that the ODE system
becomes maximally stiff, the composition is a function of state to high
accuracy (\citealt{hix1999b}: hundreds of abundances ``uniquely defined by the
thermodynamic conditions and a single measure of the weak interaction
history''), and integrating it is both expensive and pointless. KEPLER-family
codes \citep{weaver1978} and SN~Ia/CCSN hydro codes (e.g.\ \citealt{calder2007};
\citealt{seitenzahl2009}; \citealt{timmes2000} for the NSE-energy-generation
approach) therefore switch to an \textbf{NSE solve or table}: compute
$X(T, \rho, \Ye)$ from Saha, and evolve only \Ye{} through the weak reactions
evaluated on that composition. \textbf{MESA does not}: it runs the same in-situ
network at all temperatures --- ``no NSE or QSE approximation was used; the
same 204 isotope reaction network was used throughout''
(\citealt[\S 6.1]{paxton2015}, and \S 5.2 for why they avoid the switch: ad-hoc
NSE/QSE switching ``can introduce unphysical discontinuities'' and
instabilities) --- and neither does bbq (grep of the r23.05.1 and bbq trees: no
NSE control or code path). So the NSE regime below is a \emph{deployment-host}
property, not a universal one. (An earlier version of this section stated the
switch as universal stellar-code behaviour; corrected in the 2026-08-18 audit.)

That is a \emph{different mathematical object} from a network step:

\begin{center}
\begingroup\small
\begin{tabularx}{\textwidth}{@{}L{2.6cm} L{4.2cm} Y@{}}
\toprule
 & \textbf{network regime} & \textbf{NSE regime} \\
\midrule
composition & integrated ODE, path-dependent & algebraic function of
  $(T, \rho, \Ye)$ \\
free variables & $n$ abundances & \textbf{one} (\Ye) \\
what is integrated & $\dd Y/\dd t = \nu R$ & $\dd\Ye/\dd t = \sum_{\mathrm{weak}}
  (\ldots)$ evaluated at $X_{\mathrm{NSE}}$ (\citealt[\S 5]{hix1999b}: NSE
  ``requires monitoring of weak reaction activity'') \\
cost & stiff implicit solve & one Newton solve on 2 unknowns ($Y_p$, $Y_n$;
  \citealt{hix1999b}) \\
stiffness & $\ge 10^{14}$ (Tier~2 \S IV.1) & none \\
\bottomrule
\end{tabularx}
\endgroup
\end{center}

Four consequences the project should be carrying and is not --- each scoped,
after the 2026-08-18 audit, to the kind of host.

\textbf{(i) For an NSE-switching host, the emulator's competitor above the
threshold is not a network --- it is a Saha solve.} And a Saha solve is
\emph{fast}: this project's own batched solver does 27 states with no fallback
and agrees with pynucastro to $10^{-10}$~dex \resultsref{2026-07-10}. If a host
code with an NSE switch would use NSE at $\Tn > 5$ anyway, then an emulator
competing there is competing with something already cheap and already exact.
This is exactly \foreignreg{2}{R-16} / this tier's \reg{T-08}
(NSE-as-a-predictor as a baseline) arriving from the \emph{deployment} side
rather than the evaluation side, and it is a strong independent argument for
\reg{T-03}'s validity domain. For MESA the competitor at $\Tn > 5$ is the
in-situ network at its stiffest, and the argument reverts to
\S\ref{part:host}'s cost case.

\textbf{(ii) In NSE-switching hosts the handoff is discontinuous, and the
discontinuity is a physical modelling error nobody in this project has sized.}
At the switch temperature the composition jumps from ``whatever the network
integrated'' to ``exactly NSE''. Those differ, because the network was
\emph{approaching} NSE, not at it (\S\ref{sec:alpharich}'s freeze-out is the
general form of this). \citet[\S 5.2]{paxton2015} describe exactly these
discontinuities and instabilities as the reason MESA avoids the switch;
\tierone{VIII.E.6}'s ``unspecified handoff'' is the \emph{emulator
$\leftrightarrow$ solver} handoff, a different seam. The sharper statement is
that \textbf{the transition is where the label generator changes character},
and it sits inside the regime box.

\textbf{(iii) The pathology of \S\ref{part:pathology} lives on the wrong side
of it --- for NSE-switching hosts.} The Appendix-B-displaced labels are at
$\Tn \gtrsim 5$ --- i.e. in the band where an NSE-switching host would have
handed off and never called the network at all. For MESA that band is network
territory (\citealt{farmer2016} run 204 isotopes to collapse), so the ``least
deployment-relevant'' argument holds only for NSE-switching hosts, and the
benchmark-comparability tip toward \reg{T-03}'s third option is
correspondingly weaker than an earlier version of this section claimed: if the
host does not call a network there, training an emulator to reproduce buggy
labels there costs nothing scientifically and buys benchmark comparability; if
the host is MESA, it does call the network there.

It also has a sharp corollary: \textbf{the shipped test trajectories run bbq's
network at temperatures where an NSE-switching stellar code would not (MESA
would).} bbq is a bare burner with no NSE switch, so the dataset explores a
regime that one class of deployment path avoids.

\textbf{(iv) \Ye{} still evolves in NSE, and that is the \emph{only} thing that
evolves.} In the NSE regime the entire dynamical content is the weak sector
--- the same 20 channels of \S\ref{sec:yecarriers} while the Saha composition
stays in the $\Ye \gtrsim 0.46$ band (\citealt[\S 8]{paxton2015} note that
NSE-regime weak rates need large isotope pools, \citealt{juodagalvis2010}),
evaluated on a Saha composition. An emulator for that regime would be a very
different and much smaller object: a map $(T, \rho, \Ye) \to \dd\Ye/\dd t$,
three inputs and one output. Whether that is worth building separately, and
whether it is in fact the \emph{easiest and most valuable} part of the problem,
is not a question the project has asked.

Registered as \reg{T-21}.

% =====================================================================
\subsection{Self-check}
\label{sec:selfcheck-siburn}

Answer without scrolling up.

\begin{selfcheck}
% source 0.10 Q13
\item Above $\sim$5--7~GK some stellar/hydro codes replace the network with an
  NSE solve; MESA does not. Give three consequences for this project, one of
  which concerns where the label pathology lives, and say which survive when
  the host is MESA.
\end{selfcheck}
